% অনুশীলন ভান্ডার — অধ্যায় ৪
% এই ফাইলের নাম থেকেই অধ্যায় নির্ধারিত হয়।

\begin{exo}[একটি পথ বরাবর অন্তরক রূপের সমাকলন]{integrale-forme-differentielle}
\keywords{অন্তরক রূপ, সম্পূর্ণ অন্তরক, শোয়ারৎসের মানদণ্ড, রেখা সমাকলন}

দুটি অন্তরক রূপ বিবেচনা করা যাক:
\[
\omega_1 = y\,\mathrm{d}x + x\,\mathrm{d}y,
\qquad
\omega_2 = y\,\mathrm{d}x,
\]
এবং মূলবিন্দু $O(0,0)$ থেকে $M(1,1)$ বিন্দুতে যাওয়ার তিনটি পথ:
$\gamma_1$ পথে অনুভূমিক রেখাংশ $O\to(1,0)$-এর পরে উল্লম্ব রেখাংশ
$(1,0)\to M$; $\gamma_2$ পথে উল্লম্ব রেখাংশ $O\to(0,1)$-এর পরে অনুভূমিক
রেখাংশ $(0,1)\to M$; আর $\gamma_3$ হলো $O$ ও $M$-কে সরাসরি সংযোগকারী রেখাংশ।

\begin{enumerate}
\item এমন একটি $f$ অপেক্ষক দেখিয়ে প্রমাণ করুন যে $\omega_1$ সম্পূর্ণ, যেখানে $\omega_1=\mathrm{d}f$।
\item প্রতিটি রেখাংশের প্যারামিটারায়ন করে $\int_{\gamma_1}\omega_1$ ও $\int_{\gamma_2}\omega_1$ নির্ণয় করুন। হিসাব ছাড়াই ফলটি আবার পান।
\item $\omega_2$-এর উপর শোয়ারৎসের মানদণ্ড প্রয়োগ করুন। কী সিদ্ধান্ত নেওয়া যায়?
\item $\int_{\gamma_1}\omega_2$, $\int_{\gamma_2}\omega_2$ ও $\int_{\gamma_3}\omega_2$ নির্ণয় করে মন্তব্য করুন।
\end{enumerate}

\begin{indication}
অনুভূমিক রেখাংশে $\mathrm{d}y=0$ এবং উল্লম্ব রেখাংশে $\mathrm{d}x=0$;
তাই প্রতিটি সমাকলন একটি সরল সমাকলনে নেমে আসে।
\end{indication}

\begin{solution}
\textbf{প্রশ্ন ১।} $f(x,y)=xy$ উপযুক্ত, কারণ $\partial f/\partial x=y$ এবং $\partial f/\partial y=x$।

\textbf{প্রশ্ন ২।} $t\mapsto\bigl(x(t),y(t)\bigr)$ দ্বারা প্যারামিটারায়িত বক্ররেখার জন্য
\[
\int_\gamma\omega_1
= \int \left[y(t)x'(t)+x(t)y'(t)\right]\,\mathrm{d}t.
\]

$\gamma_1$ পথটি দুটি রেখাংশের সংযোগ। প্রথমটিতে $t\in[0,1]$-এর জন্য
$x(t)=t$, $y(t)=0$; তাই $\mathrm{d}x=\mathrm{d}t$ ও $\mathrm{d}y=0$।
দ্বিতীয়টিতে একইভাবে $t\in[0,1]$-এর জন্য $x(t)=1$, $y(t)=t$; তাই
$\mathrm{d}x=0$ ও $\mathrm{d}y=\mathrm{d}t$। অতএব,
\[
\begin{aligned}
I_1
= \int_{\gamma_1}\omega_1
&= \int_0^1\left(0\times 1+t\times 0\right)\,\mathrm{d}t
 + \int_0^1\left(t\times 0+1\times 1\right)\,\mathrm{d}t \\
&= 0+\int_0^1\mathrm{d}t = 1.
\end{aligned}
\]

$\gamma_2$-এর প্রথম রেখাংশে $x(t)=0$, $y(t)=t$, তারপর দ্বিতীয়টিতে
$x(t)=t$, $y(t)=1$; উভয় ক্ষেত্রেই $t\in[0,1]$। তাই
\[
\begin{aligned}
I_2
= \int_{\gamma_2}\omega_1
&= \int_0^1\left(t\times 0+0\times 1\right)\,\mathrm{d}t
 + \int_0^1\left(1\times 1+t\times 0\right)\,\mathrm{d}t \\
&= 0+\int_0^1\mathrm{d}t = 1.
\end{aligned}
\]
দুটি মান সমান। হিসাব ছাড়াই এটি অনুমান করা যেত: $f(x,y)=xy$-এর জন্য
$\omega_1=\mathrm{d}f$, তাই সমাকলনটি শুধু প্রান্তবিন্দুর উপর নির্ভর করে,
\[
\int_\gamma\omega_1=f(M)-f(O)=f(1,1)-f(0,0)=1.
\]

\textbf{প্রশ্ন ৩।} লিখি $\omega_2=A(x,y)\,\mathrm{d}x+B(x,y)\,\mathrm{d}y$।
এখানে $A(x,y)=y$ এবং $B(x,y)=0$। রূপটি সম্পূর্ণ হলে শোয়ারৎসের মানদণ্ড দিত
\[
\frac{\partial A}{\partial y}=\frac{\partial B}{\partial x}.
\]
কিন্তু $\partial A/\partial y=1$, অথচ $\partial B/\partial x=0$। মিশ্র
অবকলজ দুটি ভিন্ন; তাই $\omega_2$ সম্পূর্ণ নয়। এর সমাকলন $O$ থেকে $M$
পর্যন্ত অনুসৃত পথের উপর নির্ভর করতে পারে, যা পরের হিসাবটি নিশ্চিত করে।

\textbf{প্রশ্ন ৪।} যেহেতু $\omega_2=y\,\mathrm{d}x$, তাই $y$ শূন্য না হলে
$x$-এর পরিবর্তনই শুধু সমাকলনে অবদান রাখতে পারে। আগের প্যারামিটারায়ন থেকে
$\gamma_1$-এর জন্য পাই
\[
\begin{aligned}
J_1
=\int_{\gamma_1}\omega_2
&=\int_0^1 0\times 1\,\mathrm{d}t
  +\int_0^1 t\times 0\,\mathrm{d}t \\
&=0.
\end{aligned}
\]
$\gamma_2$-এর প্রথম রেখাংশে $x(t)=0$, তাই $\mathrm{d}x=0$; দ্বিতীয়টিতে
$x(t)=t$, $y(t)=1$ ও $\mathrm{d}x=\mathrm{d}t$। সুতরাং
\[
\begin{aligned}
J_2
=\int_{\gamma_2}\omega_2
&=\int_0^1 t\times 0\,\mathrm{d}t
  +\int_0^1 1\times 1\,\mathrm{d}t \\
&=1.
\end{aligned}
\]
শেষে, কর্ণ রেখাংশ $\gamma_3$ স্পষ্টভাবে প্যারামিটারায়িত হয়
\[
x(t)=t,\qquad y(t)=t,\qquad t\in[0,1],
\]
তাই $\mathrm{d}x=\mathrm{d}t$ ও $\mathrm{d}y=\mathrm{d}t$। ফলে
\[
J_3
=\int_{\gamma_3}\omega_2
=\int_0^1 y(t)x'(t)\,\mathrm{d}t
=\int_0^1 t\,\mathrm{d}t
=\frac12.
\]
তিনটি পথের প্রান্তবিন্দু একই, কিন্তু তিনটি ভিন্ন মান দেয়: $J_1=0$,
$J_2=1$ ও $J_3=1/2$। এটিই অসম্পূর্ণ অন্তরক রূপের বৈশিষ্ট্যপূর্ণ পথ-নির্ভরতা।
\end{solution}
\end{exo}

\begin{exo}[একই অভ্যন্তরীণ শক্তির পরিবর্তন, স্থানান্তরের তিন পদ্ধতি]{meme-delta-u-trois-transferts}
\keywords{প্রথম সূত্র, অভ্যন্তরীণ শক্তি, তাপ, যান্ত্রিক কাজ, বৈদ্যুতিক কাজ, ক্যালরিমিতি, অবস্থা অপেক্ষক}

একটি তরল, তার ক্যালরিমিটার ও নিমজ্জিত একটি ছোট রোধ মিলে স্থূলভাবে স্থির
একটি বদ্ধ ব্যবস্থা গঠন করে। পাত্রটি অনমনীয় এবং ব্যবস্থার মোট তাপধারণ ক্ষমতা
ধ্রুব বলে ধরা হয়:
\[
C=2{,}00\ \mathrm{kJ\,K^{-1}}.
\]
ব্যবস্থাটিকে $T_A=20{,}0\,^\circ\mathrm{C}$ তাপমাত্রার সাম্যাবস্থা $A$ থেকে
$T_B=25{,}0\,^\circ\mathrm{C}$ তাপমাত্রার সাম্যাবস্থা $B$-তে নিতে হবে। ধরা যাক,
\[
U(B)-U(A)=C(T_B-T_A).
\]
স্থূল গতিশক্তি ও স্থিতিশক্তির পরিবর্তন শূন্য। বাইরের দিকে সব ক্ষয় উপেক্ষা করা হয়।

একই প্রাথমিক অবস্থা $A$ থেকে নিচের তিনটি প্রক্রিয়া স্বাধীনভাবে সম্পন্ন করা হয়:
\begin{itemize}
\item[প্রক্রিয়া ১] কোনো কাজ না দিয়ে ব্যবস্থাটিকে একটি উষ্ণ উৎসের সঙ্গে তাপীয় সংস্পর্শে আনা হয়।
\item[প্রক্রিয়া ২] দেয়ালগুলি তাপরোধী। $m$ ভরের একটি বস্তু $h=5{,}00$ m উচ্চতা দিয়ে পড়লে একটি পাখা তরলটিকে নাড়ায়। শেষে পাখা ও তরল স্থির হয় এবং ভরটির হারানো সমস্ত যান্ত্রিক শক্তি ব্যবস্থায় স্থানান্তরিত হয়। $g=9{,}81\ \mathrm{m\,s^{-2}}$ নিন।
\item[প্রক্রিয়া ৩] ব্যবস্থা $Q_3=4{,}00$ kJ তাপ গ্রহণ করে এবং অবশিষ্ট শক্তি বৈদ্যুতিক পথে রোধটিতে দেওয়া হয়। গৃহীত বৈদ্যুতিক ক্ষমতা ধ্রুব, $\mathcal P=300$ W।
\end{itemize}

\begin{enumerate}
\item তিনটি প্রক্রিয়ার অভিন্ন অভ্যন্তরীণ শক্তির পরিবর্তন $\Delta U=U(B)-U(A)$ নির্ণয় করুন।
\item প্রথম দুটি প্রক্রিয়ার প্রতিটির জন্য ব্যবস্থা কর্তৃক গৃহীত তাপ $Q$ ও কাজ $W$ নির্ণয় করুন। দ্বিতীয় প্রক্রিয়ায় প্রয়োজনীয় ভর $m$ হিসাব করুন।
\item তৃতীয় প্রক্রিয়ার জন্য গৃহীত বৈদ্যুতিক কাজ ও রোধে বিদ্যুৎ দেওয়ার সময় নির্ণয় করুন।
\item তিনটি রূপান্তরের প্রাথমিক ও চূড়ান্ত অবস্থা একই। তবু তাদের $Q$ ও $W$ কেন ভিন্ন হতে পারে তা ব্যাখ্যা করুন। চূড়ান্ত অবস্থা $B$-তে ব্যবস্থাটি কি “তাপ” বা “কাজ” ধারণ করে?
\end{enumerate}

\begin{indication}
ব্যাংকারের চিহ্নরীতি মেনে $\Delta U=Q+W$ প্রয়োগ করুন। পাখার ক্ষেত্রে গৃহীত
কাজ ভরটির স্থিতিশক্তির $mgh$ হ্রাসের সমান। তারের মধ্য দিয়ে সীমানা অতিক্রমকারী
বৈদ্যুতিক শক্তিকে বৈদ্যুতিক কাজ হিসেবে গণনা করা হয়।
\end{indication}

\begin{solution}
\textbf{প্রশ্ন ১।} তাপমাত্রার পার্থক্য $T_B-T_A=5{,}0$ K। অতএব,
\[
\Delta U=C(T_B-T_A)
=2{,}00\times5{,}0
=10{,}0\ \mathrm{kJ}.
\]
এই মানটি শুধু সাম্যাবস্থা $A$ ও $B$-এর উপর নির্ভর করে।

\textbf{প্রশ্ন ২।} প্রথম প্রক্রিয়ায় কোনো কাজ গৃহীত হয় না: $W_1=0$। প্রথম সূত্র তাই দেয়
\[
Q_1=\Delta U=10{,}0\ \mathrm{kJ}.
\]
শক্তি ব্যবস্থায় প্রবেশ করে বলে তাপ ধনাত্মক।

দ্বিতীয় প্রক্রিয়ায় দেয়াল তাপরোধী, তাই $Q_2=0$। অভ্যন্তরীণ শক্তির পুরো
পরিবর্তন পাখার যান্ত্রিক কাজ থেকে আসে:
\[
W_2=\Delta U=10{,}0\ \mathrm{kJ}.
\]
আবার $W_2=mgh$, তাই
\[
m=\frac{W_2}{gh}
=\frac{10{,}0\times10^3}{9{,}81\times5{,}00}
\simeq 204\ \mathrm{kg}.
\]
মানটির বড় মাত্রা মনে করিয়ে দেয় যে সামান্য তাপমাত্রা বৃদ্ধিও ইতিমধ্যে
উল্লেখযোগ্য যান্ত্রিক শক্তির সমতুল্য।

\textbf{প্রশ্ন ৩।} প্রথম সূত্র অনুযায়ী
\[
W_3=\Delta U-Q_3=10{,}0-4{,}00=6{,}00\ \mathrm{kJ}.
\]
কাজটি ধনাত্মক: বৈদ্যুতিক সরবরাহ ব্যবস্থায় শক্তি দেয়। $W_3=\mathcal P\,\Delta t$ থেকে
\[
\Delta t=\frac{W_3}{\mathcal P}
=\frac{6{,}00\times10^3}{300}
=20{,}0\ \mathrm{s}.
\]

\textbf{প্রশ্ন ৪।} তিনটি পথ যথাক্রমে দেয়
\[
(Q_1,W_1)=(10{,}0,0)\ \mathrm{kJ},
\qquad
(Q_2,W_2)=(0,10{,}0)\ \mathrm{kJ},
\]
এবং
\[
(Q_3,W_3)=(4{,}00,6{,}00)\ \mathrm{kJ}.
\]
প্রতিটি ক্ষেত্রেই $Q+W$-এর যোগফল $10{,}0$ kJ এবং একই $\Delta U$ সৃষ্টি করে।
অভ্যন্তরীণ শক্তি একটি অবস্থা অপেক্ষক, কিন্তু তাপ ও কাজ রূপান্তরের সময় শক্তি
স্থানান্তর বর্ণনা করে। চূড়ান্ত অবস্থা $B$-তে ব্যবস্থার অভ্যন্তরীণ শক্তি আছে,
কিন্তু তাতে “তাপ” বা “কাজ” সঞ্চিত নেই।
\end{solution}
\end{exo}

\begin{exo}[ধ্রুব বহিঃচাপে সংকোচন]{compression-pression-exterieure-constante}
\seoready{true}
\keywords{চাপবলগুলির কাজ, ধ্রুব বহিঃচাপ, সংকোচন, প্রসারণ, শূন্যস্থানে প্রসারণ}

একটি গ্যাসকে ধ্রুব বহিঃচাপ $P_0$-এর অধীনে $V_A$ আয়তন থেকে $V_B<V_A$ আয়তনে সংকুচিত করা হয়।

\begin{enumerate}
\item গ্যাস কর্তৃক গৃহীত কাজ নির্ণয় করুন।
\item গ্যাস একই বহিঃচাপ $P_0$-এর বিরুদ্ধে $V_B$ থেকে $V_A$-তে ফিরে যায়। গৃহীত কাজ নির্ণয় করে সংকোচনের কাজের সঙ্গে তুলনা করুন।
\item গ্যাস যখন শূন্যস্থানে প্রসারিত হয়, তখন $V_B$ থেকে $V_A$ প্রসারণটি আবার বিবেচনা করুন। প্রাপ্ত তিনটি চিহ্ন ব্যাখ্যা করুন।
\end{enumerate}

\begin{indication}
$W=-\int P_{\mathrm{ext}}\,\mathrm{d}V$ ব্যবহার করুন। রূপান্তরটি প্রায়-স্থিতিশীল
না হলেও \emph{বহিঃচাপই} সমাকলন করতে হবে।
\end{indication}

\begin{solution}
\textbf{প্রশ্ন ১।} বহিঃচাপ ধ্রুব, তাই
\[
W_{A\to B}=-P_0(V_B-V_A)=P_0(V_A-V_B)>0.
\]
সংকোচনের সময় গ্যাস কাজ গ্রহণ করে।

\textbf{প্রশ্ন ২।} প্রত্যাবর্তনের ক্ষেত্রে
\[
W_{B\to A}=-P_0(V_A-V_B)<0.
\]
দুটি কাজ পরস্পরের বিপরীত, কারণ রূপান্তর দুটি একই ধ্রুব বহিঃচাপের বিরুদ্ধে করা হয়।

\textbf{প্রশ্ন ৩।} শূন্যস্থানে $P_{\mathrm{ext}}=0$, তাই
\[
W_{B\to A}^{\mathrm{vide}}=0.
\]
অতএব প্রসারণ মানেই স্বয়ংক্রিয়ভাবে ঋণাত্মক কাজ নয়: কোনো বহিঃবলকে বাস্তবে
পিছনে ঠেলতে হবে।
\end{solution}
\end{exo}
